%------------------------------------------------------------------------------%
\begin{question}
  Determine o conjunto solução do sistema
  \[
    \systeme{
       x + 2y -  z = 0,
      2x + 4y - 2z = 0,
       x +  y +  z = 0
    }
  \]

  \bigskip
  \begin{columns}[t]
  \column{0.3\linewidth}
    \pause
    Matriz aumentada
    \[
      \begin{amatrix}{3}
        1 & 2 & -1 & 0 \\
        2 & 4 & -2 & 0 \\
        1 & 1 & 1  & 0
      \end{amatrix}
    \]
  \column{0.4\linewidth}
    \pause
    Basta escalonar a matriz $A$
    \[
      A =
      \begin{bmatrix}
        1 & 2 & -1  \\
        2 & 4 & -2  \\
        1 & 1 &  1
      \end{bmatrix}
    \]
  \end{columns}
\end{question}

%------------------------------------------------------------------------------%
\begin{resolution}{Escalonamento}
  \begin{columns}
  \column{0.3\linewidth}
    \[\;\;\;
      \begin{bmatrix}
        1 & 2 & -1 \\
        2 & 4 & -2 \\
        1 & 1 &  1
      \end{bmatrix}
    \]
    \pause
  \column{0.7\linewidth}
    Eliminamos os elementos abaixo do primeiro pivô
    \\ \pause
    \(
      L_2 \leftarrow L_2 - 2 L_1
    \)
    \\ \pause
    \(
      L_3 \leftarrow L_3 - L_1
    \)
  \end{columns}
  \vspace{-\baselineskip}
  \[
    \pause
    L'_2
    \;=\;
    L_2 - 2 L_1
    \pause
    \;=\;
    \begin{bmatrix}
      2 & 4 & -2
    \end{bmatrix}
    - 2
    \begin{bmatrix}
      1 & 2 & -1
    \end{bmatrix}
    \pause
    \;=\;
    \begin{bmatrix}
      0 & 0  & 0
    \end{bmatrix}
  \]
  \[
    \pause
    L'_3
    \;=\;
    L_3 - L_1
    \pause
    \;=\;
    \begin{bmatrix}
      1 & 1 &  1
    \end{bmatrix}
    -
    \begin{bmatrix}
      1 & 2 & -1
    \end{bmatrix}
    \pause
    \;=\;
    \begin{bmatrix}
      0 & -1 & 2
    \end{bmatrix}
  \]
  \pause
  \[
    \begin{bmatrix}
      1 & 2  & -1 \\
      0 & 0  & 0  \\
      0 & -1 & 2
    \end{bmatrix}
  \]
\end{resolution}

%------------------------------------------------------------------------------%
\begin{resolution}{Escalonamento}
  \begin{columns}
  \column{0.3\linewidth}
    \[\;\;\;
      \begin{bmatrix}
        1 & 2  & -1 \\
        0 & 0  & 0  \\
        0 & -1 & 2
      \end{bmatrix}
    \]
    \pause
  \column{0.7\linewidth}
    Trocamos as linhas $L_2$ e $L_3$
  \end{columns}

  \pause
  \[
    \begin{bmatrix}
      1 & 2  & -1 \\
      0 & -1 & 2  \\
      0 & 0  & 0
    \end{bmatrix}
  \]
\end{resolution}

%------------------------------------------------------------------------------%
\begin{resolution}{Escalonamento}
  \begin{columns}
  \column{0.3\linewidth}
    \[\;\;\;
      \begin{bmatrix}
        1 & 2  & -1 \\
        0 & -1 & 2  \\
        0 & 0  & 0
      \end{bmatrix}
    \]
    \pause
  \column{0.7\linewidth}
    Tornando o segundo pivô igual a 1
    \\ \pause
    \(
      L_2 \leftarrow - L_2
    \)
  \end{columns}
  \vspace{-\baselineskip}
  \[
    \pause
    L'_2
    \;=\;
    - L_2
    \pause
    \;=\;
    -
    \begin{bmatrix}
      0 & -1 & 2
    \end{bmatrix}
    \pause
    \;=\;
    \begin{bmatrix}
      0 & 1 & -2
    \end{bmatrix}
  \]
  \pause
  \[
    \begin{bmatrix}
      1 & 2 & -1 \\
      0 & 1 & -2 \\
      0 & 0 & 0
    \end{bmatrix}
  \]
\end{resolution}

%------------------------------------------------------------------------------%
\begin{resolution}{Escalonamento}
  \begin{columns}
  \column{0.3\linewidth}
    \[\;\;\;
      \begin{bmatrix}
        1 & 2 & -1 \\
        0 & 1 & -2 \\
        0 & 0 & 0
      \end{bmatrix}
    \]
    \pause
  \column{0.7\linewidth}
    Eliminamos os elementos a cima do segundo pivô
    \\ \pause
    \(
      L_1 \leftarrow L_1 - 2 L_2
    \)
  \end{columns}
  \vspace{-\baselineskip}
  \[
    \pause
    L'_1
    \;=\;
    L_1 - 2 L_2
    \pause
    \;=\;
    \begin{bmatrix}
      1 & 2 & -1
    \end{bmatrix}
    - 2
    \begin{bmatrix}
      0 & 1 & -2
    \end{bmatrix}
    \pause
    \;=\;
    \begin{bmatrix}
      1 & 0 & 3
    \end{bmatrix}
  \]
  \pause
  \[
    \begin{bmatrix}
      1 & 0 & 3  \\
      0 & 1 & -2 \\
      0 & 0 & 0
    \end{bmatrix}
  \]
\end{resolution}

%------------------------------------------------------------------------------%
\begin{resolution}{Solução}
  Sistema reduzido
  \[
    \systeme{
      x   + 3z = 0,
        y - 2z = 0
    }
  \]

  \pause
  \bigskip
  A variável $z$ é livre
  \[
    z = t
    \qquad
    \forall\, t \in \R
  \]
\end{resolution}

%------------------------------------------------------------------------------%
\begin{resolution}{Solução}
  Solução
  \begin{align*}
    x & = -3t \\
    y & =  2t
  \end{align*}

  \pause
  Vetorialmente
  \[
    \mathbf{x}
    \pause
    \;=\;
    \begin{bmatrix}
      -3t \\
      2t  \\
      t
    \end{bmatrix}
    \pause
    \;=\;
    t
    \begin{bmatrix}
      -3 \\
      2  \\
      1
    \end{bmatrix}
  \]
\end{resolution}

%------------------------------------------------------------------------------%
\begin{resolution}{Conjunto Solução}
  O \emph{conjunto solução} é
  \[
    S
    =
    \left\{
    t
    \begin{bmatrix}
      -3 \\
      2  \\
      1
    \end{bmatrix}
    \st
    t\in\mathbb{R}
    \right\}
  \]

  \pause
  O conjunto solução forma uma reta que passa pela origem
\end{resolution}

%------------------------------------------------------------------------------%
\endinput
